% !TeX program = pdflatex
% Primary Article Template (ACM), sigconf option, double column.
% Content mirrors paper_wikis/paper_draft.md verbatim (2026-07-14 version).
%
% TODO before submission (cannot be filled in without your input):
%   1. Author names / affiliations / emails (Section "AUTHOR BLOCK" below) —
%      paper_draft.md itself only has placeholders ("FirstName Surname"),
%      so these were not guessed.
%   2. \acmConference / \copyrightyear / \acmYear / rights block — the
%      challenge submission portal's required rights option was not
%      specified anywhere in the repo, so \setcopyright{none} is used as a
%      neutral placeholder. Replace with whatever TAPS/the portal requires.
%   3. \acmDOI, \acmISBN — only assigned by ACM/the portal at acceptance;
%      left blank via \settopmatter{printacmref=false} + no \acmDOI call.
%
% Packages used: acmart's own load list, plus booktabs (table rules) and
% listings (verbatim anchor-format block, Section 3.1). Both are on ACM's
% TAPS-approved package whitelist. No other packages are used.
%
% Not verified to compile: no local TeX installation was available in this
% environment (pdflatex/xelatex/lualatex not found). Compile with a real
% ACM Primary Article Template checkout before submitting.

\documentclass[sigconf]{acmart}

\usepackage{booktabs}
\usepackage{listings}

\setcopyright{none}
\settopmatter{printacmref=false}
\renewcommand\footnotetextcopyrightpermission[1]{}
\pagestyle{plain}

\acmConference[RecSys Challenge 2026]{ACM RecSys 2026 Challenge: TalkPlayData Conversational Music Recommendation}{2026}{}

\lstset{
  basicstyle=\ttfamily\small,
  breaklines=true,
  columns=fullflexible,
  frame=single
}

\begin{document}

\title{A CPU-Scale Conversational Music Recommender: Fused Retrieval and Learned Ranking for the TalkPlayData Challenge}

% ---------------- AUTHOR BLOCK (placeholders — fill in before submission) ----------------
\author{FirstName Surname}
\affiliation{%
  \institution{Institution/University Name}
  \city{City}
  \state{State}
  \country{Country}
}

\author{FirstName Surname}
\affiliation{%
  \institution{Institution/University Name}
  \city{City}
  \state{State}
  \country{Country}
}

\author{FirstName Surname}
\affiliation{%
  \institution{Institution/University Name}
  \city{City}
  \state{State}
  \country{Country}
}
% -------------------------------------------------------------------------------------------

\renewcommand{\shortauthors}{Surname et al.}

\begin{abstract}
We present a conversational music recommender for the TalkPlayData Challenge
(ACM RecSys 2026) that, at each dialogue turn, ranks tracks from a 47,071-track
catalog and generates the assistant response, and that trains and serves entirely
on a single consumer desktop with no discrete GPU. Retrieval is a union of nine
inexpensive sources whose candidates are arbitrated by a gradient-boosted
learning-to-rank model over 67 features; the strongest feature is simply how many
sources agree on a candidate. The only trained neural component is a LoRA
adaptation of a 279M-parameter multilingual encoder over a role-tagged dialogue
anchor that encodes prior turns, recommendations, and per-turn user reactions. We
further identify and correct a timing misalignment and rejected-track
contamination in the dataset's goal-progress feedback signal, and show the same
signal can be reconstructed from raw dialogue when it is withheld at test time.
The system reaches a composite score of 0.49 on the blind generalization set.
\end{abstract}

\begin{CCSXML}
<ccs2012>
   <concept>
       <concept_id>10002951.10003317</concept_id>
       <concept_desc>Information systems~Recommender systems</concept_desc>
       <concept_significance>500</concept_significance>
   </concept>
   <concept>
       <concept_id>10002951.10003317.10003338</concept_id>
       <concept_desc>Information systems~Learning to rank</concept_desc>
       <concept_significance>500</concept_significance>
   </concept>
   <concept>
       <concept_id>10010147.10010178.10010179</concept_id>
       <concept_desc>Computing methodologies~Language models</concept_desc>
       <concept_significance>300</concept_significance>
   </concept>
</ccs2012>
\end{CCSXML}

\ccsdesc[500]{Information systems~Recommender systems}
\ccsdesc[500]{Information systems~Learning to rank}
\ccsdesc[300]{Computing methodologies~Language models}

\keywords{conversational recommendation, music recommendation, learning to rank, dense retrieval, LoRA fine-tuning, preference elicitation, CPU-scale systems}

\maketitle

\section{Introduction}

The TalkPlayData Challenge frames music recommendation as a multi-turn dialogue.
At each turn the system must (a) rank 20 tracks from a catalog of 47,071 so that
the single logged next track ranks high, and (b) produce the assistant's reply.
Submissions are scored by a composite metric:

\begin{quote}
Score $= 0.50 \times \text{nDCG@20} + 0.10 \times \text{CatalogDiversity} + 0.10 \times \text{LexicalDiversity} + 0.30 \times \text{JudgeNorm}$,\\
$\text{JudgeNorm} = (\text{judge\_score} - 1) / 4$.
\end{quote}

The judge is a language model that reads only the response text (personalization
and explanation quality); it never sees the track list. The blind generalization
set (``Blind B'') is harder than the labeled data: the conversation goal, the
per-turn goal-progress assessments, and the model's internal thoughts are removed,
and half of its 80 sessions are cold-start with no user profile.

Our system has three stages. \emph{Recall} runs nine retrievers per turn and merges
their outputs into one deduplicated pool (mean $\sim$3,300 candidates); no single
source exceeds $\sim$59\% gold recall at 500 candidates, but the fused pool reaches
$\sim$89\% on the blind-simulated holdout. \emph{Rescore} is a LightGBM LambdaMART model that
scores every candidate from 67 features and emits the top 20, replacing hand-tuned
score fusion. \emph{Respond} is decoupled: because the judge reads only text, the
response may describe the ranked list but never reorders it. All three stages run
on one Apple M4 machine (16~GB unified memory, no GPU), and the only trained
neural component is the LoRA encoder of Section~3.

\textbf{Contributions.} (1) A recall design that treats retrieval as \emph{set union, not
weighted fusion}: each source contributes candidates and per-source features, and
a learned ranker decides how far to trust each; multi-source agreement is the
single strongest feature. (2) A dialogue-aware retriever: a LoRA fine-tune of
\texttt{multilingual-e5-base} over a role-tagged anchor encoding prior turns,
recommendations, and per-turn reactions. (3) A treatment of the goal-progress
feedback signal that is off-by-one in time and contaminated with rejected tracks,
plus a test-time reconstruction of that signal from raw dialogue.

\section{Recall Architecture}

Recall is a union of nine retrievers (Table~\ref{tab:sources}). Each candidate carries the
per-source scores and origin flags that later become ranker features. The design
rule is to \emph{add candidates, never re-weight sources by hand}: even low-precision
sources help, because source membership feeds the agreement feature of Section~4.

\begin{table}
  \caption{The nine recall sources.}
  \label{tab:sources}
  \begin{tabular}{clll}
    \toprule
    \# & Source & Budget & Signal \\
    \midrule
    1 & BM25 (Okapi) & top-500 & lexical match on name/artist/album/tags \\
    2 & Two-tower ANN & top-2000 & fine-tuned conversational-intent match \\
    3 & Last-track NN & 100 $\times$ last 2 & ``more like what just played'' \\
    4 & Artist expansion & discographies & tracks by any retrieved/mentioned artist \\
    5 & Qwen3-0.6B NN & top-500 & second dense opinion, different space \\
    6 & CF-BPR & top-200 & collaborative signal, warm users only \\
    7 & Session-mean NN & top-100 & neighbors of the session taste centroid \\
    8 & Co-occurrence & 300/150/50 & tracks that followed recent history \\
    9 & SASRec buckets & cap 300 & next semantic-bucket expansion (optional) \\
    \bottomrule
  \end{tabular}
\end{table}

\emph{Lexical (source 1).} BM25 (via \texttt{bm25s}) indexes each track's name, artist, album,
and tag list; indexing the tags matters because conversational requests use
genre/mood/era vocabulary that only tags carry. The query concatenates the goal,
culture, the last four played tracks, and the last four dialogue turns (the
current request among them); older turns are excluded because they dilute the
current request's IDF mass.

\emph{Dense (sources 2, 5).} Two independent dense retrievers give complementary
spaces. Source 2 is our fine-tuned two-tower encoder (Section~3): the dialogue
anchor is encoded once and matched against a precomputed 47,071 $\times$ 768
L2-normalized matrix by exact brute-force matmul (top-2000); at this catalog size
an approximate index is unnecessary and only loses recall. Source 5 is the
competition-provided \texttt{Qwen3-Embedding-0.6B} metadata embedding, whose different
pretraining recovers gold tracks the two-tower misses; both cosines become ranker
features.

\emph{Neighborhood (sources 3, 7).} Query-side retrieval is complemented by expansion
around the history: the nearest neighbors of the last two played tracks (source 3)
and of the mean two-tower embedding of the last four history tracks, a recent
taste centroid (source 7). Both are cheap matmuls against the same matrix; rejected history tracks
are filtered from these seeds first (Section~3.3).

\emph{Catalog structure (sources 4, 8, 9).} Artist expansion (source 4) is a
dictionary lookup that adds the full, popularity-capped discography of every
retrieved or mentioned artist; it is the cheapest high-recall source and matters
because sessions are often artist deep-dives. Co-occurrence (source 8) is a count
table of tracks that most often followed the last one/two/three history tracks in
training sessions. Semantic-bucket expansion (source 9) is an optional auxiliary:
an RQ-VAE assigns each track a semantic ID and a small SASRec predicts the next
bucket, but next-bucket prediction has a low ceiling (hit@3 = 42.1\%), so an
eight-source configuration without it reaches parity.

\emph{Collaborative (source 6).} The competition-provided BPR matrix-factorization
embeddings add each warm user's top-200 tracks; cold sessions simply receive nothing from this source and the
pipeline degrades gracefully.

\section{Dialogue-Aware Retriever}

Off-the-shelf dense retrieval failed here: conversational queries do not embed
near track-metadata text in general-purpose spaces, so dense recall requires
fine-tuning on (dialogue, gold-track) pairs.

\subsection{Model and anchor}

The base encoder is \texttt{intfloat/multilingual-e5-base} (XLM-RoBERTa, 768-dim,
512-token window, 279M parameters), chosen for its context length (an earlier
256-token model truncated the user's request on 81\% of anchors). We fine-tune with
LoRA (rank 32, alpha 64) on the query/key/value projections: 1.8M trainable
parameters, 0.64\% of the model; full fine-tuning does not fit in 16~GB, LoRA plus
gradient checkpointing does. The query anchor is \emph{role-tagged} so the encoder sees
dialogue structure:

\begin{lstlisting}
query: [PROFILE] {age}.{country}.{gender}.{culture} [GOAL] {goal}
[T1] USER:{q1} | REC:{track-artist} | ASST:{r1} | REACTION:liked/rejected
... [NOW] USER:{current request}
\end{lstlisting}

A tokenizer-aware greedy builder always keeps a fixed core (profile, goal, latest
request) and inserts history turns most-recent-first within a 510-token budget; the
current request can never be truncated. Documents use the symmetric form
\texttt{passage: \{name\} by \{artist\} | Album: \{album\} | Tags: \{tags\} | \{year\}}, with the
E5 \texttt{query:}/\texttt{passage:} prefixes identical across training, index build, and
inference.

\subsection{Training and serving}

We train with MultipleNegativesRankingLoss (InfoNCE with in-batch negatives) plus
two explicit hard negatives per row (Section~3.3): 3 epochs, effective
batch 32, learning rate 1e-4. Training uses $\sim$46K anchor/positive pairs (plus a
7K-pair validation split), excluding the 6,000 sessions reserved for the ranker's
feature dump. At serving time we ship
a 7~MB LoRA adapter applied to the cached base encoder, and query a one-pass
47,071 $\times$ 768 track index by brute-force matmul.

\subsection{Using the goal-progress feedback signal}

Each training turn carries a \texttt{goal\_progress\_assessment} (GPA): a label recording
whether the \emph{previous} recommendation moved the listener toward their goal, i.e. a
per-recommendation accept/reject signal. It is the richest preference feedback in
the dataset and the trickiest part of the system, because the raw label cannot be
used as-is and it is removed entirely from the blind set.

\textbf{First, align it in time.} The GPA at turn $T$ is the reaction to the track
recommended at turn $T{-}1$: the system acts at one turn and the user responds at the
next. A naive reading attaches each reaction to the wrong track, so we shift every
GPA back one turn (\texttt{turn\_number - 1}). This alignment is a prerequisite for
everything below.

\textbf{Second, stop rewarding rejected tracks.} After alignment, 43.2\% of turns
carry a gold track the user \emph{rejected} (\texttt{DOES\_NOT\_MOVE\_TOWARD\_GOAL}). Training
these as ordinary positives teaches the models to rank rejected tracks highly.
The ranker \emph{drops} such turns entirely (a rejected track is not a valid
positive). The retriever additionally recycles every rejected track as a \emph{hard
negative} for the rest of that session, exactly the near-miss the encoder
should learn to push away: negatives are drawn in priority order from confirmed
session rejections, then BM25 near-misses, then same-artist distractors, then
in-batch negatives, and an accepted track is never sampled as a negative.

\textbf{Third, act on the feedback at inference.} The aligned GPA drives three
behaviors: the history-based retrievers (sources 3, 7 and the BM25 query) are
seeded only from \emph{accepted} tracks, so a streak of rejections does not fill the
pool with neighbors of refused tracks; the ranker receives features contrasting a
candidate against the accepted- and rejected-history centroids (Section~4); and
the anchor's goal slot is replaced by the most recently accepted track (``more like
the last thing that worked'').

\textbf{Finally, reconstruct it when missing.} The blind set removes GPA, so we infer
it from the dialogue: a rule-based classifier labels each user follow-up accept or
reject from keyword banks (accept: ``love it'', ``more like this''; reject: ``not
what'', ``something different''). The inferred label fills the anchor's \texttt{REACTION}
slot and drives the same seed filtering and features. On the blind-simulated
holdout the system scores as well with this proxy as with the true label, which is
what makes it viable.

\section{Learning-to-Rank Rescoring}

The rescorer is a LightGBM LambdaMART model (nDCG@20 objective) that scores every
pooled candidate and returns the top 20. It replaces hand-tuned linear fusion and
is the pivotal choice: it \emph{decouples recall from precision}, absorbs heterogeneous
features (cosines, ranks, count tables, metadata, session counters), and handles
extreme imbalance natively (one positive per $\sim$3,300 candidates) via pairwise
gradients within each turn's group. Training runs the full retrieval in dump mode
on 6,000 held-out sessions, yielding 27,166 groups and 89.7M rows; the 8.2\% of
groups whose gold track the pool missed carry no positive and are dropped,
leaving 24,951 training groups. The model has 67 features (5-fold session CV
nDCG@20 = 0.3144) and is a 0.6~MB file. The features fall into six groups (Table~\ref{tab:features}). Group B dominates:
multi-source agreement accounts for roughly 57\% of total model gain and is the
feature that transfers best when goal and GPA are withheld on the blind set.

\begin{table}
  \caption{Rescorer feature groups (67 features).}
  \label{tab:features}
  \begin{tabular}{p{0.06\linewidth}p{0.85\linewidth}}
    \toprule
    Group & Features \\
    \midrule
    A & Retrieval scores: two-tower cosine and rank; Qwen3 metadata and lyrics cosines; CLAP audio cosine; CF cosine; BM25 reciprocal rank; co-occurrence rank \\
    B & Source agreement: \texttt{n\_sources} and transforms: count of sources returning the candidate; a de-facto ensemble vote \\
    C & Track metadata: popularity percentile; years since release; request-tag overlap \\
    D & Session state (GPA): cosine to accepted- and to rejected-history centroids; artist-in-rejected-set; accepted/rejected turn counts; consecutive rejections; CF distance to history \\
    E & Intent extraction: era/genre/mood keywords parsed from the request; candidate era/genre match; album continuity; within-artist popularity and transition ranks; bucket-match counts \\
    F & Sequential ID: SASRec bucket-origin and rank (optional configuration) \\
    \bottomrule
  \end{tabular}
\end{table}

\section{Results}

Table~\ref{tab:results} lists only the changes that moved the development nDCG@20 (1,000
sessions). The two largest single jumps are non-model changes: excluding
already-played tracks from BM25 (+0.035), and replacing linear fusion with a
learned ranker, which unlocked every later feature. The goal-progress fixes of
Section~3.3 account for the final step to 0.1864. On the blind sets the same
system scores nDCG@20 = 0.3997 (Blind A) and composite 0.49 (Blind B). Two
negative results shaped the design: post-hoc editing of the ranker's top-20 to
``clean up obvious mistakes'' regressed nDCG in every trial, because logged
next-tracks often lie inside the same artist cluster a curator would prune; and
zero-shot cross-encoder and LLM rerankers underperformed the feature-based
LambdaMART while costing orders of magnitude more compute.

\begin{table}
  \caption{Development nDCG@20 after each load-bearing change.}
  \label{tab:results}
  \begin{tabular}{lr}
    \toprule
    Change & nDCG@20 \\
    \midrule
    BM25 over name + artist + album & 0.0861 \\
    + index the tag list & 0.0960 \\
    + exclude already-played tracks & 0.1313 \\
    + fine-tuned two-tower, linear fusion & 0.1418 \\
    + multi-source pool expansion & 0.1533 \\
    + LightGBM LambdaMART replaces linear fusion & 0.1609 \\
    + feature engineering, larger two-tower pool & 0.1684 \\
    + E5 role-tagged dialogue anchor & 0.1748 \\
    + GPA-aware inference and intent features & 0.1864 \\
    \bottomrule
  \end{tabular}
\end{table}

\section{Efficiency and Scalability}

Every stage is CPU-native; Table~\ref{tab:latency} (top) reports per-component latencies with all
neural blocks forced to CPU. Exact brute-force similarity over the 47,071 $\times$ 768
matrix is a single matmul, so no approximate index is needed at this catalog
scale, and every stage except the encoder passes is single-digit milliseconds.
Measured end-to-end over the 200-session evaluation set, the full recall-plus-
rescore pipeline runs at \textbf{2.17~s per session, about 0.27~s per turn} (Table~\ref{tab:latency},
bottom). Of this, the rescore stage is negligible: scoring the $\sim$3,300-candidate
pool with LambdaMART takes $\sim$3.4~ms, roughly 1\% of turn latency, so recall (the nine
sources and their two query-encoder passes) accounts for essentially all of it.
Even with every neural block pinned to the CPU (Table~\ref{tab:latency}, top), a full turn stays
under one second; dropping the frozen Qwen3-0.6B encoder in the eight-source
configuration removes the single largest cost. The trained artifact is a 24~MB
LoRA adapter and the serving process fits in under $\sim$2~GB of RAM, so the entire
train/evaluate/iterate loop, including the one LoRA fine-tune, runs on a single
16~GB desktop, which is what made rapid iteration practical.

\begin{table}
  \caption{Measured latency and footprint. Top: per component, all neural blocks
  forced to CPU (worst case), median over 20 runs. Bottom: end-to-end recall +
  rescore over the 200-session set, encoders on the Metal backend.}
  \label{tab:latency}
  \begin{tabular}{lr}
    \toprule
    Operation & Cost \\
    \midrule
    Two-tower query encode (e5-base) & 47~ms \\
    Exact ANN, 47,071 $\times$ 768, top-2000 & 2.2~ms \\
    Qwen3-0.6B query encode & 601~ms \\
    BM25 retrieve, top-500 & 0.5~ms \\
    Rescore: LambdaMART, 3,100 $\times$ 67 features & 3.4~ms \\
    Serving memory, all components loaded & 1.85~GB \\
    Trained artifact (LoRA adapter) & 24~MB \\
    Recall stage, per turn (nine sources) & $\sim$0.26~s \\
    Rescore stage, per turn & 3.4~ms \\
    End-to-end, per session (per turn) & 2.17~s (0.27~s) \\
    \bottomrule
  \end{tabular}
\end{table}

\section{Conclusion}

A conversational music recommender need not use a GPU cluster or a per-candidate
neural reranker to be competitive. Fusing many inexpensive retrievers, letting a
gradient-boosted ranker arbitrate them, adding a single LoRA-adapted dialogue
encoder, and correctly handling a goal-progress feedback signal that is off-by-one
in time, contaminated with rejected tracks, and absent at test time together reach
composite 0.49 on the blind set while training and serving on one desktop.

\begin{acks}
Insert acknowledgments here.
\end{acks}

\begin{thebibliography}{7}

\bibitem{e5}
Liang Wang, Nan Yang, Xiaolong Huang, Linjun Yang, Rangan Majumder, and Furu Wei. 2024.
Multilingual E5 Text Embeddings: A Technical Report.
\emph{arXiv preprint arXiv:2402.05672}.

\bibitem{lightgbm}
Guolin Ke, Qi Meng, Thomas Finley, Taifeng Wang, Wei Chen, Weidong Ma, Qiwei Ye, and Tie-Yan Liu. 2017.
LightGBM: A Highly Efficient Gradient Boosting Decision Tree.
In \emph{Advances in Neural Information Processing Systems (NeurIPS)}.

\bibitem{lambdamart}
Christopher J. C. Burges. 2010.
From RankNet to LambdaRank to LambdaMART: An Overview.
Microsoft Research Technical Report MSR-TR-2010-82.

\bibitem{lora}
Edward J. Hu, Yelong Shen, Phillip Wallis, Zeyuan Allen-Zhu, Yuanzhi Li, Shean Wang, Lu Wang, and Weizhu Chen. 2022.
LoRA: Low-Rank Adaptation of Large Language Models.
In \emph{International Conference on Learning Representations (ICLR)}.

\bibitem{bpr}
Steffen Rendle, Christoph Freudenthaler, Zeno Gantner, and Lars Schmidt-Thieme. 2009.
BPR: Bayesian Personalized Ranking from Implicit Feedback.
In \emph{Proceedings of the Twenty-Fifth Conference on Uncertainty in Artificial Intelligence (UAI)}.

\bibitem{sasrec}
Wang-Cheng Kang and Julian McAuley. 2018.
Self-Attentive Sequential Recommendation.
In \emph{IEEE International Conference on Data Mining (ICDM)}.

\bibitem{bm25}
Stephen Robertson and Hugo Zaragoza. 2009.
The Probabilistic Relevance Framework: BM25 and Beyond.
\emph{Foundations and Trends in Information Retrieval} 3, 4.

\end{thebibliography}

\end{document}
